\section{Relevant course concepts}
\subsection{Local generation, tokens, and RAG}
A \emph{local LLM} performs core inference on the user's own machine through a model server; in StudyGraph this server is Ollama. This reduces dependence on hosted inference, although model size and hardware still constrain quality and speed. Text is divided into tokens---often subwords rather than whole words. A transformer language model~\cite{vaswani2017attention} estimates a distribution for the next token conditional on preceding tokens and supplied context, and decoding repeats this process. Thus a generated answer is not a database lookup or a fitted numerical function.

RAG adds a retrieval step before this generation. Documents are prepared and indexed offline; at question time, relevant passages are retrieved and inserted into the model's context~\cite{lewis2020rag}. The model's parameters are not updated by a query. Retrieval can improve access to course-specific facts and provenance, but generated grounding remains contingent on passage quality and model behaviour.

\subsection{Embeddings and similarity}
An embedding model maps text to a numerical vector. Vectors learned from usage patterns can place semantically related text near one another~\cite{mikolov2013efficient}; a query vector can therefore be compared with stored document vectors by a vector index. This differs from token generation: embeddings support retrieval, while the chat LLM produces the answer.

The course's skip-gram demonstration made this learning process concrete. New tokens initially had essentially random representations; when they entered training, their loss rose, after which gradient updates moved their vectors and their loss fell. A three-dimensional PCA view showed these movements, but PCA is only a projection for visualisation, not the embedding model itself. This was a course example, not an embedding strategy we implemented or optimised. In our MVP we call a pretrained local embedding model and do not claim to have evaluated semantic-embedding quality systematically.

\subsection{Supervised regression and evaluation}
For one predictor, linear regression fits $\hat y=ax+b$ by estimating $a$ and $b$ from training pairs $(x_i,y_i)$, conventionally by minimising squared residuals. Held-out data assess generalisation rather than fit to rows used for training. We report mean absolute error,
\begin{equation}
\mathrm{MAE}=\frac{1}{n}\sum_{i=1}^{n}\left|y_i-\hat y_i\right|,
\end{equation}
which has the target's units and weights each absolute residual equally. Unlike the LLM arm, the fitted model learns task-specific coefficients. Scikit-learn supplies the implementation used here~\cite{pedregosa2011scikit}.
